% GENERATED FILE. Edit canonical lessons, then run npm run generate:books.
% canonical-source-hash: fnv1a64:c86758e7be492d80
% canonical-lessons: ZH-C419-jieji, ZH-C419-duihuan, ZH-C419-waibi, ZH-C419-guoji, ZH-C419-guonei, ZH-R419-review-414-419

\chapter{Arrivals}
\label{ch:zh-arrivals}
\hlchaptermodality{419}

\begin{chapteropening}
\textbf{By the end of this chapter:} \emph{I can meet someone off a plane, change money, and say a flight is international or domestic.}

Complete the last lesson of chapter 419: I can meet someone off a plane, change money, and say a flight is international or domestic.
\end{chapteropening}

\section[jiējī]{jiējī — to meet someone off a plane}

\label{lesson:ZH-C419-jieji}



\begin{quote}
Before the new one: say the Mandarin for to deal with the paperwork, then the Mandarin for over the weight limit.
\end{quote}

\subsection*{You'll want to know: jiējī}
\emph{jiējī} — \textquotedblleft{}to meet someone off a plane\textquotedblright{}.

\begin{quote}
\emph{míngtiān wǒ zuò chūzūchē qù jīchǎng jiējī} — Tomorrow I'm taking a taxi to the airport to meet someone.
\end{quote}

Its characters come later; for now it is a word for the ear.

\begin{cousinweb}
\emph{jiējī} is built from \emph{jiē} 'to meet' and \zh{机} 'a plane', both inherited from Middle Chinese, \emph{tsjep} and \emph{kj+j} (Guangyun, in Baxter's transcription), which glosses the first 'to join' and 'to meet'.
\end{cousinweb}

\begin{grammarlens}[title={What you've built}]
One more word for getting around.
\end{grammarlens}

\subsection*{Your turn}
\begin{itemize}
\raggedright
  \item \emph{Say it:} \emph{jiējī}
  \item \emph{Say it:} \emph{jiējī}, once more
  \item \emph{Say it:} \emph{míngtiān wǒ zuò chūzūchē qù jīchǎng jiējī} — Tomorrow I'm taking a taxi to the airport to meet someone.
  \item \emph{Recall:} say the Mandarin for to deal with the paperwork, then the Mandarin for over the weight limit, then say \emph{jiējī} again
\end{itemize}

\begin{tcolorbox}[breakable,colback=teal!4,colframe=teal!35!black,title={Before you move on}]
What does jiējī mean? (\textquotedblleft{}To meet someone off a plane\textquotedblright{}.) Say it once more.
\end{tcolorbox}
\section[duìhuàn]{duìhuàn — to exchange money}

\label{lesson:ZH-C419-duihuan}



\begin{quote}
Before the new one: say the Mandarin for over the weight limit, then the Mandarin for to meet someone off a plane.
\end{quote}

\subsection*{You'll want to know: duìhuàn}
\emph{duìhuàn} — \textquotedblleft{}to exchange money\textquotedblright{}.

\begin{quote}
\emph{wǒ zài jīchǎng duìhuàn le yìxiē qián} — I changed some money at the airport.
\end{quote}

Its characters come later; for now it is a word for the ear.

\begin{cousinweb}
\emph{duìhuàn} is built from \emph{duì} and \emph{huàn} 'to exchange', both inherited from Middle Chinese, \emph{dwajH} and \emph{hwanH} (Guangyun, in Baxter's transcription), which glosses the first as one of the eight trigrams and 'to please', and the second 'to exchange'; the sense 'to exchange' for the first is not given there.
\end{cousinweb}

\begin{grammarlens}[title={What you've built}]
One more word for getting around.
\end{grammarlens}

\subsection*{Your turn}
\begin{itemize}
\raggedright
  \item \emph{Say it:} \emph{duìhuàn}
  \item \emph{Say it:} \emph{duìhuàn}, once more
  \item \emph{Say it:} \emph{wǒ zài jīchǎng duìhuàn le yìxiē qián} — I changed some money at the airport.
  \item \emph{Recall:} say the Mandarin for over the weight limit, then the Mandarin for to meet someone off a plane, then say \emph{duìhuàn} again
\end{itemize}

\begin{tcolorbox}[breakable,colback=teal!4,colframe=teal!35!black,title={Before you move on}]
What does duìhuàn mean? (\textquotedblleft{}To exchange money\textquotedblright{}.) Say it once more.
\end{tcolorbox}
\section[wàibì]{wàibì — foreign currency}

\label{lesson:ZH-C419-waibi}



\begin{quote}
Before the new one: say the Mandarin for to meet someone off a plane, then the Mandarin for to exchange money.
\end{quote}

\subsection*{You'll want to know: wàibì}
\emph{wàibì} — \textquotedblleft{}foreign currency\textquotedblright{}.

\begin{quote}
\emph{zhèr kěyǐ huàn wàibì} — You can change foreign money here.
\end{quote}

Its characters come later; for now it is a word for the ear.

\begin{cousinweb}
\emph{wàibì} is built from \zh{外} 'outside' and \emph{bì} 'money', both inherited from Middle Chinese, \emph{ngwajH} and \emph{bjiejH} (Guangyun, in Baxter's transcription), which glosses the second as gifts of silk.
\end{cousinweb}

\begin{grammarlens}[title={What you've built}]
One more word for getting around.
\end{grammarlens}

\subsection*{Your turn}
\begin{itemize}
\raggedright
  \item \emph{Say it:} \emph{wàibì}
  \item \emph{Say it:} \emph{wàibì}, once more
  \item \emph{Say it:} \emph{zhèr kěyǐ huàn wàibì} — You can change foreign money here.
  \item \emph{Recall:} say the Mandarin for to meet someone off a plane, then the Mandarin for to exchange money, then say \emph{wàibì} again
\end{itemize}

\begin{tcolorbox}[breakable,colback=teal!4,colframe=teal!35!black,title={Before you move on}]
What does wàibì mean? (\textquotedblleft{}Foreign currency\textquotedblright{}.) Say it once more.
\end{tcolorbox}
\section[guójì]{guójì — international}

\label{lesson:ZH-C419-guoji}



\begin{quote}
Before the new one: say the Mandarin for to exchange money, then the Mandarin for foreign currency.
\end{quote}

\subsection*{You'll want to know: guójì}
\emph{guójì} — \textquotedblleft{}international\textquotedblright{}.

\begin{quote}
\emph{qù guójì hángzhànlóu zěnme zǒu} — How do I get to the international terminal?
\end{quote}

Its characters come later; for now it is a word for the ear.

\begin{cousinweb}
\emph{guójì} is built from \zh{国} 'a country' and \emph{jì} 'an edge, a meeting', both inherited from Middle Chinese, \emph{kwok} and \emph{tsjejH} (Guangyun, in Baxter's transcription), which glosses the second 'a border' and 'to meet'.
\end{cousinweb}

\begin{grammarlens}[title={What you've built}]
One more word for getting around.
\end{grammarlens}

\subsection*{Your turn}
\begin{itemize}
\raggedright
  \item \emph{Say it:} \emph{guójì}
  \item \emph{Say it:} \emph{guójì}, once more
  \item \emph{Say it:} \emph{qù guójì hángzhànlóu zěnme zǒu} — How do I get to the international terminal?
  \item \emph{Recall:} say the Mandarin for to exchange money, then the Mandarin for foreign currency, then say \emph{guójì} again
\end{itemize}

\begin{tcolorbox}[breakable,colback=teal!4,colframe=teal!35!black,title={Before you move on}]
What does guójì mean? (\textquotedblleft{}International\textquotedblright{}.) Say it once more.
\end{tcolorbox}
\section[guónèi]{guónèi — domestic (inside the country)}

\label{lesson:ZH-C419-guonei}



\begin{quote}
Before the new one: say the Mandarin for foreign currency, then the Mandarin for international.
\end{quote}

\subsection*{You'll want to know: guónèi}
\emph{guónèi} — \textquotedblleft{}domestic (inside the country)\textquotedblright{}.

\begin{quote}
\emph{guónèi hángbān hěn piányi} — Domestic flights are cheap.
\end{quote}

Its characters come later; for now it is a word for the ear.

\begin{cousinweb}
\emph{guónèi} is built from \zh{国} 'a country' and \emph{nèi} 'inside', both inherited from Middle Chinese, \emph{kwok} and \emph{nwojH} (Guangyun, in Baxter's transcription), which glosses the second 'to go in'.
\end{cousinweb}

\begin{grammarlens}[title={What you've built}]
One more word for getting around.
\end{grammarlens}

\subsection*{Your turn}
\begin{itemize}
\raggedright
  \item \emph{Say it:} \emph{guónèi}
  \item \emph{Say it:} \emph{guónèi}, once more
  \item \emph{Say it:} \emph{guónèi hángbān hěn piányi} — Domestic flights are cheap.
  \item \emph{Recall:} say the Mandarin for foreign currency, then the Mandarin for international, then say \emph{guónèi} again
\end{itemize}

\begin{tcolorbox}[breakable,colback=teal!4,colframe=teal!35!black,title={Before you move on}]
What does guónèi mean? (\textquotedblleft{}Domestic (inside the country)\textquotedblright{}.) Say it once more.
\end{tcolorbox}
\section[(dialogue)]{Looking back: from In the Cabin to Arrivals}

\label{lesson:ZH-R419-review-414-419}



\begin{quote}
Say the Mandarin for a flight attendant and next to the window.
\end{quote}

\subsection*{Your turn}
\begin{itemize}
\raggedright
  \item \emph{Say it:} a flight attendant, next to the window, economy class, to get airsick, a seat belt
  \item \emph{Say it:} an airport terminal building, to check in a suitcase, a boarding pass, the security check at an airport, a departure gate
  \item \emph{Say it:} to be held up, the measure word for a rainstorm, to miss a flight, to change planes on the way, flying (of an aircraft)
  \item \emph{Say it:} customs at the border, to enter a country, to leave a country, to declare goods at customs, duty-free
  \item \emph{Say it:} to check in for a flight, the formalities, to be grounded, to deal with the paperwork, over the weight limit
  \item \emph{Say it:} to meet someone off a plane, to exchange money, foreign currency, international, inside the country
  \item \emph{Read it:} \zh{场} \zh{停} — each one aloud, then a word that holds it
\end{itemize}

\begin{tcolorbox}[breakable,colback=teal!4,colframe=teal!35!black,title={Before you move on}]
A flight attendant? (\emph{chéngwùyuán}.) Next to the window? (\emph{kàochuāng}.) Economy class? (\emph{jīngjìcāng}.) To get airsick? (\emph{yùnjī}.) A seat belt? (\emph{ānquándài}.) An airport terminal building? (\emph{hángzhànlóu}.) To check in a suitcase? (\emph{tuōyùn}.) A boarding pass? (\emph{dēngjīpái}.) The security check at an airport? (\textbf{\zh{安检}}, \emph{ānjiǎn}.) A departure gate? (\emph{dēngjīkǒu}.) To be held up? (\emph{yánwù}.) The measure word for a rainstorm? (\textbf{\zh{场}}, \emph{chǎng}.) To miss a flight? (\emph{wùjī}.) To change planes on the way? (\emph{zhuǎnjī}.) Flying (of an aircraft)? (\emph{fēixíng}.) Customs at the border? (\emph{hǎiguān}.) To enter a country? (\emph{rùjìng}.) To leave a country? (\emph{chūjìng}.) To declare goods at customs? (\emph{shēnbào}.) Duty-free? (\emph{miǎnshuì}.) To check in for a flight? (\emph{zhíjī}.) The formalities? (\emph{shǒuxù}.) To be grounded? (\textbf{\zh{停飞}}, \emph{tíngfēi}.) To deal with the paperwork? (\emph{bànlǐ}.) Over the weight limit? (\emph{chāozhòng}.) To meet someone off a plane? (\emph{jiējī}.) To exchange money? (\emph{duìhuàn}.) Foreign currency? (\emph{wàibì}.) International? (\emph{guójì}.) Inside the country? (\emph{guónèi}.)
\end{tcolorbox}
